\documentclass[11pt]{article}

\usepackage[T1]{fontenc}
\usepackage[utf8]{inputenc}
\usepackage{lmodern}
\usepackage{geometry}
\usepackage{microtype}
\usepackage{amsmath,amssymb,amsthm,mathtools}
\usepackage{enumitem}
\usepackage[colorlinks=true,linkcolor=blue,urlcolor=blue]{hyperref}
\geometry{margin=1in}
\setlist[enumerate]{leftmargin=*,itemsep=3pt,topsep=3pt}

\newtheorem{theorem}{Theorem}[section]
\newtheorem{proposition}[theorem]{Proposition}
\newtheorem{lemma}[theorem]{Lemma}
\newtheorem{corollary}[theorem]{Corollary}
\theoremstyle{remark}
\newtheorem{remark}[theorem]{Remark}

\newcommand{\R}{\mathbb{R}}
\newcommand{\T}{\mathbb{T}}
\newcommand{\sS}{\mathcal{S}}
\newcommand{\sA}{\mathcal{A}}
\newcommand{\sE}{\mathcal{E}}
\newcommand{\sF}{\mathcal{F}}
\newcommand{\ii}{\mathrm{i}}
\newcommand{\norm}[1]{\left\lVert #1\right\rVert}
\newcommand{\abs}[1]{\left\lvert #1\right\rvert}
\newcommand{\inner}[2]{\left\langle #1,#2\right\rangle}
\DeclareMathOperator{\supp}{supp}
\DeclareMathOperator{\esssup}{ess\,sup}

\title{Fourier-Shell Symbol and Off-Shell Damping\\
for the Swift--Hohenberg Operator}
\author{M3 + E57 \quad | \quad Manuscript 1\\[2pt]
\small Self-contained mathematical proof}
\date{October 26, 2026}

\begin{document}
\maketitle

\begin{abstract}
We show that the Swift--Hohenberg shell operator
$\sA=(\Delta+k_\star^2)^2$ has the exact Fourier multiplier
$p(\xi)=(\abs{\xi}^2-k_\star^2)^2$.  For the dissipative linear equation
$\partial_tu=-a\sA u$, $a>0$, its $L^2$ solution is given by the multiplier
$e^{-atp(\xi)}$, and all frequencies separated from the preferred shell
$\abs{\xi}=k_\star$ decay at a uniform exponential rate.  We prove
self-adjointness, an energy identity, whole-space strong decay, and the
periodic-domain analogue.  We distinguish shell selection from nonlinear
pattern selection and localized bound-state existence.
\end{abstract}

\section{Definitions and main theorem}

Fix a dimension $d\ge 1$, a preferred wavenumber $k_\star>0$, and a
dissipation coefficient $a>0$.  We use the unitary Fourier transform
\begin{equation}\label{eq:fourier-convention}
\widehat u(\xi)=\sF u(\xi)
=(2\pi)^{-d/2}\int_{\R^d}e^{-\ii x\cdot\xi}u(x)\,dx,
\end{equation}
first for $u\in\sS(\R^d)$ and then on $L^2(\R^d)$ by Plancherel.
Define
\begin{equation}\label{eq:A-def}
\sA=(\Delta+k_\star^2)^2,
\qquad
p(\xi)=(\abs{\xi}^2-k_\star^2)^2.
\end{equation}

\begin{theorem}[M3 + E57: exact spectral shell damping]\label{thm:main}
For every $u_0\in L^2(\R^d)$, the initial-value problem
\begin{equation}\label{eq:flow}
\partial_tu=-a(\Delta+k_\star^2)^2u,
\qquad u(0)=u_0,
\end{equation}
has a unique $L^2$ mild solution $u\in C([0,\infty);L^2(\R^d))$.
The following statements hold:
\begin{enumerate}[label=\textup{(\roman*)}]
\item The natural $L^2$ realization of $\sA$ is a nonnegative
self-adjoint operator with domain $H^4(\R^d)$, and
\begin{equation}\label{eq:symbol}
\boxed{\widehat{\sA u}(\xi)
=(\abs{\xi}^2-k_\star^2)^2\widehat u(\xi).}
\end{equation}
\item The solution satisfies, for almost every $\xi\in\R^d$,
\begin{equation}\label{eq:semigroup}
\boxed{\widehat u(t,\xi)
=e^{-at(\abs{\xi}^2-k_\star^2)^2}\widehat u_0(\xi).}
\end{equation}
\item For $\delta>0$, let $P_\delta$ be the orthogonal Fourier projection
with multiplier $\mathbf 1_{\Omega_\delta}$, where
\begin{equation}\label{eq:omega}
\Omega_\delta
=\{\xi\in\R^d:\abs{\abs{\xi}^2-k_\star^2}\ge\delta\}.
\end{equation}
Then
\begin{equation}\label{eq:off-shell}
\boxed{\norm{P_\delta u(t)}_{L^2}
\le e^{-a\delta^2t}\norm{P_\delta u_0}_{L^2}}
\qquad (t\ge0).
\end{equation}
\item Every $L^2$ solution converges strongly to zero:
\begin{equation}\label{eq:strong-decay}
\lim_{t\to\infty}\norm{u(t)}_{L^2}=0.
\end{equation}
Nevertheless, for every finite $t\ge0$, the evolution operator
$S(t)u_0=u(t)$ satisfies
\begin{equation}\label{eq:operator-norm}
\norm{S(t)}_{L^2\to L^2}=1.
\end{equation}
Thus decay on the whole space has no uniform exponential rate.
\end{enumerate}
\end{theorem}

\section{Fourier symbol, positivity, and the operator domain}

\begin{proof}[Proof of Theorem~\ref{thm:main}(i)]
For $u\in\sS(\R^d)$ and every coordinate $j$,
\begin{equation}
\widehat{\partial_j u}(\xi)=\ii\xi_j\widehat u(\xi),
\qquad
\widehat{\partial_j^2u}(\xi)=-\xi_j^2\widehat u(\xi).
\end{equation}
Summing over $j$ gives
\begin{equation}
\widehat{\Delta u}(\xi)=-\abs{\xi}^2\widehat u(\xi).
\end{equation}
Hence
\begin{align}
\widehat{(\Delta+k_\star^2)u}(\xi)
&=(k_\star^2-\abs{\xi}^2)\widehat u(\xi),\\
\widehat{(\Delta+k_\star^2)^2u}(\xi)
&=(k_\star^2-\abs{\xi}^2)^2\widehat u(\xi)\\
&=(\abs{\xi}^2-k_\star^2)^2\widehat u(\xi),
\end{align}
which establishes the asserted Fourier symbol on Schwartz functions.

Under the unitary Fourier transform, the maximal multiplication operator
$M_p:f\mapsto pf$ is self-adjoint because $p$ is real-valued, with domain
\begin{equation}\label{eq:operator-domain}
D(M_p)=\{f\in L^2(\R^d):pf\in L^2(\R^d)\}.
\end{equation}
Since $p(\xi)\ge0$, this operator is nonnegative.  Moreover,
\begin{equation}\label{eq:equivalence}
1+p(\xi)^2\asymp(1+\abs{\xi}^2)^4
\qquad\text{for all }\xi\in\R^d,
\end{equation}
where $f\asymp g$ means that $c g\le f\le Cg$ for constants
$0<c\le C<\infty$ depending only on $k_\star$.
To see this, the ratio of the two continuous, strictly positive
functions tends to $1$ as $\abs{\xi}\to\infty$ and is bounded above and
away from zero on any compact set.  Therefore
\begin{equation}
\sF^{-1}D(M_p)=H^4(\R^d).
\end{equation}
The operator $\sA=\sF^{-1}M_p\sF$ is thus self-adjoint and
nonnegative on $H^4$; density extends the Fourier identity from
$\sS$ to this domain.

Finally, for $u\in H^4(\R^d)$, Plancherel yields
\begin{align}
\inner{u}{\sA u}_{L^2}
&=\int_{\R^d}p(\xi)\abs{\widehat u(\xi)}^2\,d\xi\\
&=\norm{(\Delta+k_\star^2)u}_{L^2}^2\ge0.
\end{align}
The associated quadratic form has the larger form domain $H^2(\R^d)$.
\end{proof}

\section{The exact dissipative evolution}

\begin{proof}[Proof of Theorem~\ref{thm:main}(ii)]
For $t\ge0$ define the Fourier multiplier
\begin{equation}\label{eq:S}
S(t)=\sF^{-1}M_{e^{-atp}}\sF.
\end{equation}
Because $0<e^{-atp(\xi)}\le1$,
\begin{equation}
\norm{S(t)u_0}_2\le\norm{u_0}_2.
\end{equation}
The identity $S(t+s)=S(t)S(s)$ follows directly from multiplication
of the scalar exponential multipliers.  Since
$e^{-atp(\xi)}\to1$ as $t\downarrow0$, dominated convergence and
Plancherel imply
\begin{equation}
\norm{S(t)u_0-u_0}_2\longrightarrow0.
\end{equation}
Thus $(S(t))_{t\ge0}$ is a strongly continuous contraction semigroup.
Its generator is $-a\sA$, which gives existence and uniqueness of the
mild solution to~\eqref{eq:flow}.  Explicitly,
\begin{equation}
\widehat u(t,\xi)=e^{-atp(\xi)}\widehat u_0(\xi)
\end{equation}
for almost every $\xi$.

For $t>0$, the functions $p(\xi)^j e^{-atp(\xi)}$ are bounded for
every nonnegative integer $j$.  Accordingly the solution belongs to the
operator domain at positive time and satisfies the differential equation
strongly there.  In particular, high spatial frequencies are smoothed.
\end{proof}

\begin{remark}[Meaning of a Fourier mode]
For general $L^2$ data, Fourier transforms are equivalence classes
defined almost everywhere.  The multiplier identity gives an exact
decay factor for almost every frequency; a plane wave of exactly fixed
frequency is a generalized eigenfunction, not an element of
$L^2(\R^d)$.
\end{remark}

\section{Uniform damping away from the shell}

\begin{proof}[Proof of Theorem~\ref{thm:main}(iii)]
On $\Omega_\delta$ as defined in~\eqref{eq:omega},
\begin{equation}
p(\xi)=(\abs{\xi}^2-k_\star^2)^2\ge\delta^2.
\end{equation}
Therefore, for $t\ge0$,
\begin{equation}
e^{-2atp(\xi)}\le e^{-2a\delta^2t}
\qquad(\xi\in\Omega_\delta).
\end{equation}
By Plancherel and~\eqref{eq:semigroup},
\begin{align}
\norm{P_\delta u(t)}_2^2
&=\int_{\Omega_\delta}
 e^{-2atp(\xi)}\abs{\widehat u_0(\xi)}^2\,d\xi\\
&\le e^{-2a\delta^2t}
 \int_{\Omega_\delta}\abs{\widehat u_0(\xi)}^2\,d\xi\\
&=e^{-2a\delta^2t}\norm{P_\delta u_0}_2^2.
\end{align}
Taking square roots proves~\eqref{eq:off-shell}.
\end{proof}

\begin{corollary}[Exact preferred shell and near-shell expansion]
The damping exponent
\begin{equation}\label{eq:lambda}
\lambda(\xi)=-a(\abs{\xi}^2-k_\star^2)^2
\end{equation}
is nonpositive everywhere, with $\lambda(\xi)=0$ exactly when
$\abs{\xi}=k_\star$.  Writing $\abs{\xi}=k_\star+h$ yields
\begin{equation}
p(\xi)=(2k_\star h+h^2)^2
=4k_\star^2h^2+4k_\star h^3+h^4.
\end{equation}
Thus the leading damping rate near the selected shell is quadratic in
the radial deviation $h$.
\end{corollary}

\section{Whole-space decay, but no uniform exponential gap}

\begin{proof}[Proof of Theorem~\ref{thm:main}(iv)]
The zero set of the symbol is
\begin{equation}\label{eq:shell}
\Sigma_\star=\{\xi\in\R^d:\abs{\xi}=k_\star\}.
\end{equation}
For $d=1$ it consists of two points, and for $d\ge2$ it is a
codimension-one sphere.  In each case it has $d$-dimensional Lebesgue
measure zero.  Hence $p(\xi)>0$ for almost every $\xi$ and
\begin{equation}
e^{-2atp(\xi)}\longrightarrow0
\qquad(t\to\infty)
\end{equation}
for almost every $\xi$.  As
$0\le e^{-2atp(\xi)}\abs{\widehat u_0(\xi)}^2
\le\abs{\widehat u_0(\xi)}^2$ and $\widehat u_0\in L^2$, dominated
convergence gives
\begin{equation}
\norm{S(t)u_0}_2^2
=\int_{\R^d}e^{-2atp(\xi)}\abs{\widehat u_0(\xi)}^2\,d\xi
\longrightarrow0.
\end{equation}

On the other hand, the norm of a bounded Fourier multiplier on $L^2$
is the essential supremum of its absolute value, so
\begin{equation}
\norm{S(t)}_{L^2\to L^2}
=\esssup_{\xi\in\R^d}e^{-atp(\xi)}=1.
\end{equation}
Indeed, every open neighborhood of the nonempty sphere $\Sigma_\star$
has positive Lebesgue measure and contains values of $p$ arbitrarily
close to zero.  No bound
$\norm{S(t)}_{L^2\to L^2}\le C e^{-ct}$ with fixed $c>0$ can therefore
hold for all $t\ge0$.  This proves the claimed distinction between
strong decay of each $L^2$ solution and uniform decay of the semigroup.
\end{proof}

\begin{corollary}[No nonzero exact-shell $L^2$ state]
If $u\in L^2(\R^d)$ and
$\supp\widehat u\subseteq\Sigma_\star$ in the essential-support sense,
then $u=0$ almost everywhere.
\end{corollary}
\begin{proof}
An $L^2$ Fourier transform supported on a Lebesgue-null set must
vanish almost everywhere.  Apply Plancherel.
\end{proof}

\section{Energy-dissipation identities}

\begin{proposition}[Monotone quadratic energy]\label{prop:energy}
Let $u(t)=S(t)u_0$ with $u_0\in L^2(\R^d)$.  Define, for $t>0$,
\begin{equation}\label{eq:energy}
\sE[u(t)]=\frac12\norm{(\Delta+k_\star^2)u(t)}_2^2.
\end{equation}
Then, for every $t>0$,
\begin{align}
\frac{d}{dt}\sE[u(t)]
&=-a\norm{(\Delta+k_\star^2)^2u(t)}_2^2\le0,
\label{eq:energy-identity}\\
\frac12\frac{d}{dt}\norm{u(t)}_2^2
&=-a\norm{(\Delta+k_\star^2)u(t)}_2^2\le0.
\label{eq:mass-identity}
\end{align}
\end{proposition}
\begin{proof}
Write $L=\Delta+k_\star^2$, so $\sA=L^2$.  At every positive
time the semigroup regularizes sufficiently for the following
differentiations and integrations by parts.  Since $L$ and $\sA$ are
self-adjoint and $u_t=-a\sA u$,
\begin{align}
\frac{d}{dt}\frac12\norm{Lu}_2^2
&=\operatorname{Re}\inner{Lu}{Lu_t}_{L^2}\\
&=\operatorname{Re}\inner{L^2u}{u_t}_{L^2}\\
&=-a\norm{\sA u}_2^2.
\end{align}
Similarly,
\begin{align}
\frac12\frac{d}{dt}\norm{u}_2^2
&=\operatorname{Re}\inner{u}{u_t}_{L^2}\\
&=-a\inner{u}{\sA u}_{L^2}\\
&=-a\norm{Lu}_2^2.
\end{align}
This proves both identities.  If $u_0\in H^4$, they are also valid
at $t=0$; otherwise they hold for $t>0$ and have the corresponding
integrated formulations.
\end{proof}

\begin{remark}[Gradient-flow convention]
Using the \emph{real} Hilbert-space inner product
$(u,v)_{\R}=\operatorname{Re}\inner{u}{v}_{L^2}$, the gradient of
$\sE[u]=\frac12\norm{Lu}_2^2$ is $\nabla_{\R}\sE[u]=\sA u$.
Thus $u_t=-a\nabla_{\R}\sE[u]$ is precisely a negative gradient flow.
This avoids possible factors of two from differing conventions for
complex Wirtinger derivatives.
\end{remark}

\section{Periodic-domain counterpart}

\begin{proposition}[Discrete Fourier shells on a flat torus]
Consider a flat torus $\T_L^d=(\R/L\mathbb Z)^d$.  Its Fourier
frequencies are $\xi_n=2\pi n/L$, $n\in\mathbb Z^d$.  Let $\Pi_0$
denote orthogonal projection onto the resonant span
\begin{equation}
\{e^{\ii\xi_n\cdot x}:\abs{\xi_n}=k_\star\},
\end{equation}
with $\Pi_0=0$ if that span is trivial.  Then for a constant $g_L>0$,
\begin{equation}\label{eq:periodic-gap}
\norm{u(t)-\Pi_0u_0}_2
\le e^{-ag_Lt}\norm{u_0-\Pi_0u_0}_2,
\end{equation}
where
\begin{equation}
g_L=\min\{(\abs{\xi_n}^2-k_\star^2)^2:
 n\in\mathbb Z^d,\ \abs{\xi_n}\ne k_\star\}>0.
\end{equation}
In particular, $u(t)\to\Pi_0u_0$ exponentially in $L^2$.
\end{proposition}
\begin{proof}
Fourier coefficients solve the decoupled equations
\begin{equation}
\widehat u_n(t)
=e^{-at(\abs{\xi_n}^2-k_\star^2)^2}\widehat u_n(0).
\end{equation}
There are finitely many lattice frequencies in each bounded ball, and
$(\abs{\xi_n}^2-k_\star^2)^2\to\infty$ as $\abs n\to\infty$.
Consequently, the nonzero values of the discrete symbol have a strictly
positive minimum $g_L$.  Parseval's identity gives
\begin{align}
\norm{u(t)-\Pi_0u_0}_2^2
&=\sum_{\abs{\xi_n}\ne k_\star}
 e^{-2at(\abs{\xi_n}^2-k_\star^2)^2}
 \abs{\widehat u_n(0)}^2\\
&\le e^{-2ag_Lt}\norm{u_0-\Pi_0u_0}_2^2,
\end{align}
where the coefficient normalization is chosen so that Parseval holds.
Taking square roots proves the claim.
\end{proof}

\section{Linear control parameter and limits of the claim}

\begin{proposition}[Linear Swift--Hohenberg threshold]
For the parameterized linear equation
\begin{equation}\label{eq:r-flow}
\partial_tu=ru-a(\Delta+k_\star^2)^2u,
\qquad r\in\R,
\end{equation}
the Fourier growth exponent is
\begin{equation}\label{eq:r-growth}
\lambda_r(\xi)=r-a(\abs{\xi}^2-k_\star^2)^2.
\end{equation}
On the continuum, the spectral bound is $r$.  Thus the linear
instability threshold in this centered parameterization is $r_c=0$;
when $r>0$, the growing frequency band is exactly
\begin{equation}\label{eq:unstable-band}
\abs{\abs{\xi}^2-k_\star^2}<\sqrt{r/a}.
\end{equation}
On $\T_L^d$, the first growing mode instead occurs when
$r>a\min_{n\in\mathbb Z^d}(\abs{\xi_n}^2-k_\star^2)^2$.
\end{proposition}
\begin{proof}
Applying $\sF$ to~\eqref{eq:r-flow} gives
$\partial_t\widehat u=\lambda_r(\xi)\widehat u$.
Since $p\ge0$ and its minimum is zero on the continuum shell,
$\sup_\xi\lambda_r(\xi)=r$; the strict positivity condition
$\lambda_r(\xi)>0$ is equivalent to~\eqref{eq:unstable-band}.
For a torus, take the maximum over its discrete Fourier frequencies.
\end{proof}

\begin{remark}[Scope: linear damping versus nonlinear persistence]
The exact coefficientwise decay formula~\eqref{eq:semigroup} concerns
the \emph{linear dissipative} evolution with the negative sign in front
of the shell operator.  Nonlinearities such as $-g\abs{u}^2u$ couple
Fourier frequencies and do not obey independent multiplier decay.
The shell symbol alone proves neither nonzero whole-space $L^2$ bound
states nor the existence, uniqueness, or nonlinear stability of
localized patterns.  In particular, changing the sign to
$+a(\Delta+k_\star^2)^2$ reverses the damping and creates
high-frequency growth.
\end{remark}

\section*{Conclusion and classification}

The principal identities are
\begin{equation*}
\boxed{\begin{gathered}
\widehat{(\Delta+k_\star^2)^2u}(\xi)
=(\abs{\xi}^2-k_\star^2)^2\widehat u(\xi),\\[4pt]
\widehat u(t,\xi)
=e^{-at(\abs{\xi}^2-k_\star^2)^2}\widehat u_0(\xi),\\[4pt]
\norm{P_\delta u(t)}_2
\le e^{-a\delta^2t}\norm{P_\delta u_0}_2.
\end{gathered}}
\end{equation*}
Accordingly, \textbf{M3 (shell Fourier symbol) is proved}, and
\textbf{E57 (off-shell spectral damping) is proved} for the stated
negative-generator linear flow.  Extensions to nonlinear persistence,
physical mode confinement, or an independently specified WCT equation
require additional hypotheses and independent theorems.

\end{document}